\subsection{The joint \texorpdfstring{$III^*$}{III*} and II caps}
\begin{statement}{P3F.47}{Proved}
On the minimal regular rational elliptic surface
\[
 y^2=x^3+tx-t,\qquad G=(1,1),
\]
use $M=\mathcal O(aG-O)$ and the actual suspension translation
$-bG$, for $(a,b)=(2,1)$ or $(1,2)$. The fibres are
$III^*$ at infinity, II at zero, and $I_1$ at $t_c=-27/4$.
Choose the simple whole zero of the line map at that $I_1$.
The following are actual simultaneous root-line cap domains:
\[
\begin{array}{c|c|c}
(a,b)&m_*\text{ at }III^*&m_0\text{ at II}\\ \hline
(2,1)&4l+3, l\geq0&6k+1, k\geq0,\quad\text{or }6k, k\geq1\\
(1,2)&4l, l\geq1&6k+1, k\geq0,\quad\text{or }6k, k\geq1
\end{array}                                            \tag{F47.1}
\]
Retain signed $n_*,n_0$ with $\gcd(m_i,n_i)=1$, including
$n_0=0$ at an order-one ordinary II cap. Put
\[
                         D=m_0n_*+m_*n_0.              \tag{F47.2}
\]
For the constructed joint family,
$\pi_1(Y)=\pi_1(X)=\mathbb Z/|D|$, with $\mathbb Z$ for $D=0$.
For $D\ne0$ the orbit groups are
\[
 H_0(Y)=H_5(Y)=\mathbb Z,\qquad
 H_1(Y)=H_3(Y)=\mathbb Z/|D|,\qquad H_2(Y)=H_4(Y)=0.
\]
For $D=0$ every $H_j(Y)$, $0\leq j\leq5$, is $\mathbb Z$.
The actual sphere tuples are exactly $D=\pm1$, for every
permitted cusp fan. The condition forces $\gcd(m_*,m_0)=1$;
in particular the good $III^*$ and good II roots cannot occur
together in a sphere tuple. This construction and computation
retain all line characters and collar coefficients.
\end{statement}
\begin{proof}
\par\medskip\noindent\textbf{The source, its actual roots, and the line maps.}\enspace
The source and its primitive height-$1/2$ section were derived
in P3F.37 and P3F.41. Its displayed discriminant has precisely
the three stated fibres. The integral source calculation below
also verifies the marked height and the full lattice.
For $h=t_{-bG}$, the exact section formula gives
$h^*M\otimes M^{-1}=\mathcal O(F)$ in both cases $ab=2$.
Its section with simple zero at $t_c$ supplies the actual forward
line map $B_0$, and $B_\lambda=\lambda B_0$ fixes its scalar
parameter and direction.

At infinity put $q=1/t$, $x_0=q^2x$, $y_0=q^3y$.
Then
\[
 y_0^2=x_0^3+q^3x_0-q^5,\qquad G=(q^2,q^3).
\]
For $m_*=4l+3$, set $q=z^{m_*}$ and divide $x_0,y_0$
by $z^{2(3l+2)},z^{3(3l+2)}$. The exact result is
\[
 Y^2=X^3+zX-z^{2l+3},\qquad
 G'=(z^{2l+2},z^{3l+3}).                            \tag{F47.3}
\]
Resolving its ordinary double point gives the reduced III
root of P3F.37. The marked $G'$ meets its nonidentity component,
whereas $2G'$ meets the identity component. For $m_*=4l$,
divide instead by $z^{6l},z^{9l}$ to obtain
\[
 Y^2=X^3+X-z^{2l},\qquad G'=(z^{2l},z^{3l}).         \tag{F47.4}
\]
Its good central fibre has deck $(X,Y)\mapsto(-X,-iY)$,
and $G'_0=(0,0)$ is the nonzero fixed two-torsion point.
Thus $bG'_0=O$ for the good pair, while the central line
$\mathcal O(G'_0-O)$ retains exact order two.

At zero use $q=t$. For $m_0=6k+1$ and $m_0=6k$, respectively,
put $q=z^{m_0}$, $x=z^{2k}X$, $y=z^{3k}Y$.
The actual equations are
\[
 \begin{aligned}
 Y^2&=X^3+z^{2k+1}X-z &&(m_0=6k+1),\\
 Y^2&=X^3+z^{2k}X-1   &&(m_0=6k).
 \end{aligned}                                         \tag{F47.5}
\]
The first has a reduced II fibre; the second has good reduction
with order-six central deck. The marked section is
$(X,Y)=(z^{-2k},z^{-3k})$ on the affine overlap. In the actual
smooth chart at $O$, its parameter is $-X/Y=-z^k$.
Hence its good central limit, and the limit of $bG'$, is $O$.
The formal group law gives $w(aG')=-az^k+O(z^{2k})$.
This proves the central section data used in the line comparison,
including the further contact of nonminimal root orders.

Write $N_i=\mathcal O(aG'_i-O_i)$ on either root. At a
reduced root the actual resolved fibre-product comparison
descends the pulled line through each contracted exceptional
curve, since its degree there is zero. Its difference from
$N_i$ is a vertical divisor with zero component degrees.
The exact III or II intersection matrix makes it an integral
whole fibre. Dividing its base monomial supplies the line
isomorphism and retains the corresponding deck character.

At the good starred root the minimal fourth deck has weights
$(3,1)$ at each fixed section. The common smooth toric refinement
has rays
\[
 (1,0),(3,1),(2,1),(1,1),(1,2),(1,3),(0,1).
\]
For its root blowup $\beta$ and resolved quotient map $\pi$,
the Cartier support function of the section is
$\max(v_w-3v_z,0)$. Consequently
\[
 \beta^*N_* =\pi^*M+D_{aG'}-D_O,\qquad
 D_P=\sum_v\min(v_w,3v_z)E_{P,v}.                    \tag{F47.6}
\]
The five exceptional coefficients are $(3,2,1,1,1)$.
The actual translation by $bG'$ at the minimal fourth root
has form $w\mapsto w\,u(w,z)+z^3h(z^4)$, with $u$ a unit.
It preserves the ideals $(w,z^j),(z,w^j)$ for $j\leq3$,
hence the normalized common blowup and (F47.6).
The further degree-$l$ root pulls this diagram back, including
its exceptional coefficients and the specified line map.

At the good II root, the minimal sixth deck has one point of
stabilizer six, one orbit of two points of stabilizer three,
and one orbit of three points of stabilizer two. Blow them up.
Their tangent weights are $(1,1)$. The quotient central curve
and exceptional images have squares $-1,-6,-3,-2$ and fibre
\[
                          6C+D_6+2D_3+3D_2.
\]
Contract $C$, then $D_2$, then $D_3$. The remaining curve
has square zero and the actual invariant Weierstrass map has
the cusp parametrization $(u^2,u^3)$ of the original II fibre.
Both $aG'$ and $O$ meet the exceptional divisor over $O$ at
finite $w/z$; all three contraction centres are at its infinity.
Each section has exceptional coefficient one. They cancel in
the exact divisor identity
\[
                  \beta^*N_0=\pi^*\alpha^*M.           \tag{F47.7}
\]
Here $\alpha$ denotes the stated contractions. Translation by
$bG'$ fixes the original blowup centres and preserves the
diagram and this equality. Pulling the diagram through the
additional degree-$k$ root proves the same assertion for
every good order in (F47.1). These are the actual line
correspondences of P3F.41 and P3F.43, now with the marked
$aG'$ and $bG'$ of this source explicitly verified.

\par\medskip\noindent\textbf{Commuting norms and the simultaneous family.}\enspace
At either root, choose the actual $D_i$-invariant division
section $R_i$, $m_iR_i=bG'_i$. At a reduced root multiplication
by $m_i$ is invertible on both the additive central identity
group and its component group. At a good root take the unique
analytic inverse branch with central value $O$.
The extending source action $t_{-n_iR_i}D_i$ has $m_i$th
power $t_{-n_i bG'_i}$. Resolve its addition graph as in
P3F.41--43; the source maps and their inverses extend on the
minimal regular root.
Lift $t_{-R_i}$ first to the specified line $N_i$.
Its exact deck commutator is a retained character $\chi_i$.
Writing $B_0=\beta_i(q_i)F_i^{m_i}$ gives the invariant
base unit $\beta_i$ and the norm factor
\[
 (F_i^{n_i}D_i)^{m_i}/B_0^{n_i}
    =\chi_i^{n_im_i(m_i-1)/2}\beta_i(q_i)^{-n_i}.
\]
Its invariant analytic $m_i$th root gives the actual commuting
line maps $A_{i,0}$ with $A_{i,0}^{m_i}=B_0^{n_i}$.
This retains the order-two central line in (F47.6) and every
scalar character in (F47.7).

Put $L=\operatorname{lcm}(m_*,m_0)$, $\lambda=r^L$, and
$A_{i,r}=r^{Ln_i/m_i}A_{i,0}$. For $n_iu_i+m_iv_i=1$,
the literal primitive generator
\[
 H_i=A_{i,r}^{u_i}B_\lambda^{v_i},\qquad
 H_i^{m_i}=B_\lambda,\quad H_i^{n_i}=A_{i,r}
\]
contains the positive scalar power $r^{L/m_i}$. Its norm shell
is uniformly contracting on smaller compact source tubes for
small $r$. Thus $N_i^\times/\langle H_i\rangle$ is a smooth
proper holomorphic cap. A unit line rotation equal to $H_i^j$
forces $j=0$ by contraction, so the delta circle is free there.
For outward $h_i=H_i^{-1}$ the actual cut marking is
\[
                 \alpha_i\longmapsto-n_i h_i,\qquad
                 \gamma\longmapsto m_i h_i.          \tag{F47.8}
\]
The holomorphic annular comparison is obtained by the actual
constant-factor degree-zero line uniformization, removing its
period two-cocycle on the annulus and then its holomorphic
Lie-bundle cocycle. Its positive cut is $A_{i,r}$ itself.
For $c=\log\lambda/(2\pi i)$ its line variation is exactly
$n_i(c-c_0)s_i\delta/m_i$; elliptic and delta periods are
retained in this equation.

For one positive $r$-loop set $k_i=Ln_i/m_i$. The ordinary
inverse closing has the global multiplier
\[
                 t^{k_0}(t-t_c)^{-(k_*+k_0)}.         \tag{F47.9}
\]
Its orders at zero and infinity are $k_0,k_*$, and its ratios
to the prescribed local multipliers are extending units.
The cap charts close by the identity in these root coordinates.
At the selected cusp, use the specified $B_\lambda$ iterates
to replace the remaining integral power by its bounded source
translation and invertible line unit. They and their inverses
extend in the actual unimodular Mumford charts. This proves
the simultaneous proper punctured-disc family and fixes its
parameter closing. For $D=\pm1$, $L=m_*m_0$, so its cusp
gamma coefficient is exactly the signed integer $D$.

\par\medskip\noindent\textbf{The source cone and all integral boundary maps.}\enspace
Use the basis $(O,F,C_0,\ldots,C_6,G)$.
The $E_7$ graph has edges
$(0,1),(1,2),(2,3),(3,4),(4,5),(2,6)$ and Cartan matrix $A$.
Its highest-root vector is $v=(2,3,4,3,2,1,2)^t$.
The Gram matrix has $O^2=G^2=-1$, $OF=GF=1$,
root block $-A$, and $G C_5=1$, with every other unspecified
pairing zero. Its determinant is $-1$.
Put $R$ equal to the inverse-Cartan column at $C_5$,
$\varphi=G-O-F+R$, and $R_j=R$ for odd $j$, zero for even $j$.
Then the exact section is
\[
 P_j=O+\left(1+\frac{j^2/2-2+c_j}{2}\right)F
                 +j\varphi-R_j,\quad
 c_j=\begin{cases}3/2&j\text{ odd},\\0&j\text{ even}.
 \end{cases}                                         \tag{F47.10}
\]
Its columns are integral, its square is $-1$, and
$\varphi^2=-1/2$. Translation sends $O$ to $G$, $G$ to $P_2$,
and reverses the two long affine $E_7$ arms. These data define
the full integral isometry $T_+$ in the certificate;
the actual source map is $T=T_+^{-b}$.
Set
\[
 L_*=(F-\sum v_jC_j,C_0,\ldots,C_6),\qquad L_0=F,
 \qquad L=(L_*,L_0).
\]
The source exterior complex is
\[
 C_4=\mathbb Z, C_3=\mathbb Z^2, C_2=\mathbb Z^{10},
 \ C_1=\mathbb Z^9, C_0=\mathbb Z,\qquad
 d_4=(1,1)^t,\quad d_2=L^tB.                         \tag{F47.11}
\]
For each boundary the complex is
$C_3=\mathbb Z$, $C_2=C_1=\mathbb Z^{r_i}$, $C_0=\mathbb Z$,
with $d_2=L_i^tBL_i$. Its inclusions are $L_i$, the component
embedding, and the $i$th coordinate in degrees two, one, three.
The actual mapping-torus differential is
\[
 d_{{\rm cone},j}
       =\begin{pmatrix}d_j&T_{j-1}-1\\0&-d_{j-1}\end{pmatrix},
\]
where the action on component coordinates is the inverse
transpose of their permutation, and the degree-three and
degree-four actions are identity. The program exports every
one of these integer matrices, every inclusion, and every
Smith representative and projection. The even free source-pair
argument in P3F.46 verifies that these are the actual derived
maps without an omitted odd homotopy term.

The exterior $W$ has $H_1=H_3=\mathbb Z^2$, $H_4=\mathbb Z$,
and $H_2=\mathbb Z^2$ for $b=1$, or $\mathbb Z/2\oplus\mathbb Z^2$
for $b=2$. In the physical degree-one marking,
$\alpha_*\mapsto\alpha$, $\alpha_0\mapsto-\alpha$,
and both gammas map to $\gamma$.
In degree three the corresponding markings are
$N_*\mapsto N$, $N_0\mapsto-N$, and both $F\gamma$ map to $Z$.
The certificate records the exact unimodular changes from its
raw Smith bases to these marked bases.

At the starred boundary choose $r_*$ with $F=2r_*$ for the
reduced root, or $F=r_*$ for the good root. At the II boundary
choose $F=r_0$. Actual specialization gives $r_i\mapsto C_i$.
Keep $s_i\mapsto\kappa_i C_i$ without restricting either
integer. In domain order $(r_*,s_*,r_0,s_0)$ the complete
degree-two matrix for $(a,b)=(2,1)$ is
\[
 \begin{pmatrix}
 0&-1&0&0\\0&-2&0&1\\
 -1&-\kappa_*&0&0\\0&0&-1&-\kappa_0
 \end{pmatrix}.                                      \tag{F47.12}
\]
For $(a,b)=(1,2)$ prepend the boundary and exterior torsion
coordinates, both of order two. The full mixed-group map is
\[
 \begin{pmatrix}
 1&1&1&1&0\\0&0&1&0&0\\0&0&1&0&1\\
 0&-1&-\kappa_*&0&0\\0&0&0&-1&-\kappa_0
 \end{pmatrix}.                                      \tag{F47.13}
\]
The first row is read modulo two. The determinants of these
free representatives are $1,-1$, respectively, independently
of the two collar coefficients. Each respects the specified
order-two relation and is an isomorphism of the actual groups.
The original unreduced integer representatives of the torsion
row are also retained in the certificate.

At a good starred cap the full boundary elliptic torsion maps
isomorphically to its order-two cap torsion. This is the
degree-one transport through the smooth root. The elliptic
lattice quotients by $D-1$ and $D^u-1$ agree since $u$ is odd.
Keep arbitrary torsion path coefficients $c,d$ in the images
of $\alpha_*,\gamma$. With column order
$(\eta,\alpha_*,\gamma_*,\alpha_0,\gamma_0)$, the full map is
\[
 D_1=\begin{pmatrix}
 0&1&0&-1&0\\0&0&1&0&1\\-1&-c&-d&0&0\\
 0&n_*&-m_*&0&0\\0&0&0&n_0&-m_0
 \end{pmatrix},                                      \tag{F47.14}
\]
whose third target coordinate has order two. Eliminating that
exact torsion isomorphism leaves, for both source pairs,
\[
 D_1^{\rm free}=D_3=
 \begin{pmatrix}
 1&0&-1&0\\0&1&0&1\\n_*&-m_*&0&0\\0&0&n_0&-m_0
 \end{pmatrix}.                                      \tag{F47.15}
\]
For degree three the cap row is $(-n_i,m_i)$ on $(N_i,F\gamma_i)$.
The actual invariant central component trace is primitive for
the reduced roots; smooth-root fibre transport has degree one
for the good roots. These facts prove both entries, with no
source three-class left over in a cap that retracts to its
two-dimensional central source and suspension circle.
The two determinants in (F47.15) are $-D$.
Finally $D_4=(1,-1)$ and
$D_0=\left(\begin{smallmatrix}1&1\\-1&0\\0&-1\end{smallmatrix}\right)$.
These give every nonzero gluing degree.

\par\medskip\noindent\textbf{The full group and the sphere conclusion.}\enspace
Let $v$ be the retained $I_1$ vanishing cycle and let $A$ be
the $III^*$ monodromy on the elliptic lattice. Then
$A^2=-I$, $\operatorname{tr}A=0$. The product relation with
the II monodromy, whose trace is one, gives
$|\langle v,Av\rangle|=1$: substitute the transvection
$T_v=I+v\otimes\langle v,\cdot\rangle$ into
$\operatorname{tr}(A^{-1}T_v^{-1})$.
Thus the vanishing relation and its conjugate by the actual
starred base loop kill an integral elliptic basis in the source
exterior. The source section then identifies its remaining
fundamental group with the base annulus group $\mathbb Z$.
The suspension adds the commuting gamma generator, giving
$\pi_1(W)=\mathbb Z\alpha\oplus\mathbb Z\gamma$.

At each good cap the same boundary lattice maps onto the full
elliptic normal subgroup, so it dies in the gluing. Reduced
caps have simply connected central source support. The exact
remaining presentation is
\[
 \langle h_*,h_0\mid h_*^{m_*}=h_0^{m_0},\quad
                         h_*^{n_*}=h_0^{-n_0}\rangle .
\]
Choose $n_*u+m_*v=1$. Substitution gives
$h_*=h_0^{-n_0u+m_0v}$, so the group is cyclic.
Its primitive relation and determinant give precisely
$\mathbb Z/|D|$. The same substitution proves that the first
three Smith factors of (F47.15) are one, also when $D=0$.
Together with the isomorphisms (F47.12)--(F47.13) and the
degree-zero and degree-four maps, the full integral sequence
gives the groups in the statement.
The symbolic identities and 1681 further exact Smith cases
are executed in \path{code/part3_root/iii_star_ii_joint.py}.

For $D=\pm1$ the orbit is simply connected with the integral
homology of $S^5$. Its standard diffeomorphism type, the
primitive fixed-sphere Hopf link, and the Gysin calculation of
P3K.22 give the actual standard oriented $S^6$ total space.
For other $D$ the displayed fundamental group excludes a
sphere. The finite common flop chain of P3F.35 carries this
topology criterion through every allowed cusp fan. All line
and annular maps used to obtain it remain the maps in
(F47.6)--(F47.9).
\end{proof}
